

\section{L'impasse actuelle : Un enseignement fossile sans moyens mais avec prétentions}
space{0.3cm}
oindent

L'enseignement du grec ancien en France traverse une crise qui ne relève plus de la simple érosion
cyclique, mais d'une mutation structurelle profonde, qualifiée par certains observateurs et acteurs
du terrain d'« impasse actuelle ». Le point I.3.\cite{I-3-3-ref3} de l'analyse critique sur
l'enseignement des langues anciennes, formulé comme « L'impasse actuelle : un enseignement fossile
qui a perdu ses moyens (heures) mais gardé ses prétentions (analyser Platon sans savoir le lire) »,
cristallise une tension insoutenable au cœur du système éducatif français. Cette formulation
lapidaire dévoile une schizophrénie institutionnelle : d'un côté, une ambition académique inchangée,
héritière des Humanités classiques du XIXe siècle, qui vise la formation d'esprits capables
d'accéder directement aux textes fondateurs de la pensée occidentale ; de l'autre, une réalité
matérielle et pédagogique dégradée, marquée par la raréfaction des horaires, la discontinuité des
apprentissages et l'effondrement des compétences linguistiques réelles.

Le terme « fossile » est ici d'une pertinence redoutable. En géologie, un fossile est une structure
minéralisée qui a conservé la forme exacte du vivant tout en ayant perdu sa substance organique. De
la même manière, l'enseignement du grec ancien a conservé sa morphologie externe — les programmes
prestigieux, les exercices canoniques de la version et du commentaire, les concours d'excellence
comme l'Agrégation — mais a été vidé de sa substance vitale : le temps d'exposition à la langue
nécessaire pour en acquérir la maîtrise. Ce rapport se propose d'explorer cette « fossilisation » en
disséquant les mécanismes historiques, politiques et didactiques qui ont conduit à cette situation
paradoxale où l'institution exige d'analyser Platon sans donner les moyens de le lire.

Nous analyserons d'abord la « perte des moyens », non comme une fatalité, mais comme le résultat de
politiques éducatives successives qui, de la réforme du collège de 2016 à celle du lycée de 2019,
ont transformé le grec en une variable d'ajustement budgétaire. Nous examinerons ensuite la «
persistance des prétentions », en confrontant les programmes officiels et les exigences des jurys de
concours à la réalité du niveau des élèves et des étudiants. Enfin, nous plongerons au cœur du «
simulacre », cette pédagogie de la compensation qui, par le recours systématique à la traduction et
au texte bilingue, tente de masquer l'analphabétisme disciplinaire tout en maintenant l'illusion
d'une transmission philologique.



\subsection{La dépossession des moyens : chronique d'une érosion programmée}

La première partie de l'équation de l'impasse réside dans la perte objective des moyens alloués à
l'enseignement du grec ancien. Cette dépossession n'est pas un accident conjoncturel mais le fruit
d'une lente érosion, accélérée brutalement par les réformes de la dernière décennie. Elle se
manifeste par une réduction drastique des volumes horaires, une précarisation du statut de la
discipline et une rupture de la continuité pédagogique.



\subsection{L'effondrement des horaires et la concurrence interne (2016-2024)}

L'histoire récente de l'enseignement des Langues et Cultures de l'Antiquité (LCA) est marquée par
une rupture majeure : la réforme du collège de 2016. Avant cette date, bien que les effectifs
fussent faibles (environ 1,10 \% de la population scolarisée dans le public en 2010 pour le grec au
lycée 1), l'enseignement bénéficiait encore de créneaux horaires identifiés. La réforme de 2016 a
été vécue comme un véritable « cataclysme » par la communauté enseignante, introduisant une
précarité structurelle inédite.\cite{I-3-3-ref2}

Le mécanisme de cette dépossession repose sur l'intégration des LCA dans les Enseignements Pratiques
Interdisciplinaires (EPI) et le financement des heures de langue sur la « marge d'autonomie » des
établissements. Concrètement, cela signifie que les heures de grec ne sont plus garanties
nationalement de la même manière, mais doivent être prélevées sur une enveloppe globale d'heures
supplémentaires, mettant la discipline en concurrence directe avec d'autres dispositifs jugés
prioritaires, comme les dédoublements de classes en sciences ou l'aide
personnalisée.\cite{I-3-3-ref3} Les conséquences chiffrées de cette réforme sont alarmantes :

\begin{itemize}\item Les horaires de LCA ont été réduits en moyenne de \textbf{70 \%} suite à la mise en œuvre de la réforme.\cite{I-3-3-ref3}\item L'enseignement du grec ancien a purement et simplement disparu de \textbf{40 \%} des collèges qui le proposaient auparavant.\cite{I-3-3-ref3}\item Dans de nombreux établissements, le maintien d'une offre de grec s'est fait au prix d'un « lissage » des horaires, réduisant souvent l'enseignement à une ou deux heures hebdomadaires, voire à un couplage artificiel avec le latin sur un créneau unique.\cite{I-3-3-ref2}\end{itemize}L'assouplissement de la réforme en 2017-2018 par le ministre Jean-Michel Blanquer, s'il a permis de
rétablir l'option facultative, n'a pas résolu le problème fondamental du financement. Le grec reste
tributaire de la volonté du chef d'établissement et de la pression des familles, ce qui accentue les
inégalités territoriales et sociales.\cite{I-3-3-ref2}



\subsection{La réforme du lycée et la marginalisation de la spécialité}

La réforme du lycée de 2019 a parachevé cette dynamique de marginalisation. En supprimant les séries
(L, S, ES) pour instaurer un système de spécialités, la réforme a théoriquement valorisé les LCA en
créant la spécialité « Littérature, Langues et Cultures de l'Antiquité » (LLCA). Cependant, la
réalité des choix d'orientation et des contraintes d'emploi du temps a produit l'effet inverse.

La spécialité LLCA est devenue une « spécialité rare », souvent fermée faute d'effectifs suffisants
ou regroupée en réseaux d'établissements, ce qui décourage les élèves. Quant à l'enseignement
optionnel, il survit difficilement face à la charge de travail imposée par les trois spécialités de
Première et les deux de Terminale. Les programmes officiels indiquent que l'enseignement optionnel
et de spécialité sont régis par des arrêtés de 2019, mais la mise en œuvre locale révèle une grande
disparité.\cite{I-3-3-ref6} Les élèves doivent souvent choisir entre le grec et des options
stratégiques pour leur orientation post-bac (Mathématiques expertes, etc.), ce qui conduit à une «
hémorragie » des effectifs entre la classe de Troisième et la Seconde, puis tout au long du cycle
terminal.\cite{I-3-3-ref1}

Le tableau suivant illustre l'évolution des contraintes structurelles pesant sur l'enseignement du
grec :

\begin{table}[h!]\small\centering\begin{tabular}{| p{2cm} | p{2.2cm} | p{2.2cm} | p{2.2cm} |}\hline\textbf{Période / Réforme} & \textbf{Statut du Grec Ancien} & \textbf{Mécanisme de Financement} & \textbf{Impact sur les Moyens (Heures)} \\ \hline\textbf{Avant 2016} & Option facultative traditionnelle & Dotation Horaire Globale (DHG) fléchée & Horaires sanctuarisés (souvent 3h), mais concurrence avec les sections européennes. \\ \hline\textbf{Réforme Collège 2016} & EPI / Enseignement de complément & Marge d'autonomie de l'établissement & Effondrement des horaires (-70\%), disparition dans 40\% des collèges offreurs.\cite{I-3-3-ref3} Concurrence mortifère avec les groupes de niveau. \\ \hline\textbf{Ajustement 2017} & Option facultative rétablie & Marge d'autonomie maintenue & Stabilisation précaire. Les heures perdues sont rarement récupérées intégralement.\cite{I-3-3-ref2} \\ \hline\textbf{Réforme Lycée 2019} & Spécialité LLCA \+ Option & DHG (Spécialité) / Marge (Option) & Création d'une niche d'excellence (Spécialité) à très faibles effectifs. L'option devient résiduelle face à la pression des coefficients du Bac.\cite{I-3-3-ref6} \\ \hline\end{tabular}\end{table}

\subsection{La pénurie d'enseignants : le cercle vicieux}

La perte des moyens est également humaine. Le rapport de l'IGEN de 2011 alertait déjà sur une «
grave pénurie d’enseignants en lettres classiques ».\cite{I-3-3-ref1} Cette pénurie s'est aggravée,
créant un cercle vicieux : la baisse du nombre d'heures et d'élèves entraîne une réduction des
postes aux concours (CAPES, Agrégation), ce qui diminue l'attractivité de la filière universitaire.
En retour, le manque d'enseignants certifiés conduit les rectorats à ne pas ouvrir ou à fermer des
sections de grec, faute de personnel pour les assurer.\cite{I-3-3-ref1}

Cette situation aboutit à des bricolages pédagogiques, comme l'enseignement du grec par des
professeurs de Lettres Modernes non spécialistes, ou le regroupement de plusieurs niveaux (Seconde,
Première, Terminale) dans un même créneau horaire d'une ou deux heures.\cite{I-3-3-ref1} Dans ces
conditions, toute progression linguistique sérieuse devient illusoire.



\subsection{La persistance des prétentions : une ambition déconnectée}

Face à cette érosion matérielle, on pourrait s'attendre à un ajustement des objectifs pédagogiques.
Or, c'est précisément l'inverse qui se produit : l'institution scolaire a « gardé ses prétentions ».
Les programmes, les examens et la rhétorique institutionnelle continuent de projeter sur les élèves
une exigence philologique que le système ne permet plus de satisfaire.



\subsection{Des programmes de lycée d'une exigence paradoxale}

L'examen des programmes de lycée pour les années 2024-2026 révèle un décalage saisissant. En
Terminale, le programme limitatif de la spécialité LLCA impose l'étude conjointe de
\textit{L'Assemblée des femmes} d'Aristophane et de \textit{La Servante écarlate} de Margaret
Atwood.\cite{I-3-3-ref8} Si la mise en regard thématique (la place des femmes, la cité,
l'utopie/dystopie) est intellectuellement riche, l'exigence formelle reste celle de l'accès au texte
grec.

Les textes officiels rappellent que « la lecture des œuvres et des textes majeurs de la littérature
gréco-latine, situés dans leur contexte, constitue le socle de l'apprentissage ».\cite{I-3-3-ref9}
Il est demandé aux élèves de se confronter à la langue d'Aristophane, connue pour sa richesse
lexicale (le vocabulaire de la comédie, les néologismes, les registres de langue) et ses
particularités dialectales (l'attique familier, les parodies de dialectes). Comment des élèves
disposant, dans le meilleur des cas, de trois ou quatre heures par semaine (souvent amputées), et
ayant subi une scolarité au collège en mode « dégradé », peuvent-ils « lire » Aristophane?

De même, la référence constante à Platon comme horizon indépassable de la formation — « analyser
Platon » — témoigne de cette persistance. Platon requiert une maîtrise syntaxique fine : l'usage des
particules (mén, dé, ara, dè), la syntaxe des modes (l'optatif oblique, le subjonctif
d'éventualité), la structure de la période. Les rapports de jury soulignent régulièrement que ces
compétences ne sont pas acquises, même par les candidats aux concours
d'enseignement.\cite{I-3-3-ref10}



\subsection{Le déni de réalité des concours de recrutement}

Le maintien des prétentions est particulièrement visible dans les rapports des jurys du CAPES et de
l'Agrégation de Lettres Classiques. Ces documents constituent une archive précieuse du déni
institutionnel. Le rapport du jury du CAPES 2020 note par exemple une moyenne de l'épreuve de langue
(latin-grec) inférieure à 10/20 (9,85 pour les Lettres Classiques) et pointe des « identifications
erronées en grec » et une « maîtrise insuffisante de la langue ».\cite{I-3-3-ref11}

Le jury déplore que les candidats, futurs professeurs, maîtrisent mal la proposition relative en
grec ou confondent des fonctions grammaticales de base.\cite{I-3-3-ref11} Pourtant, l'épreuve reine
reste la version (traduction du grec vers le français), exercice d'excellence qui suppose une
compétence que le système universitaire peine lui-même à produire, faute d'un vivier d'étudiants
solidement formés dans le secondaire.

Ainsi, le système exige des enseignants qu'ils transmettent un savoir (la lecture fluide du grec)
qu'ils n'ont eux-mêmes pas toujours eu le temps de consolider, créant une chaîne de transmission
fragilisée où la prétention d'excellence se heurte au principe de réalité.



\subsection{« analyser platon » : l'emblème de l'ambition mal placée}

La formule du point I.3.3, « analyser Platon sans savoir le lire », pointe spécifiquement le fossé
entre l'ambition philosophique et la compétence philologique.

Analyser Platon, dans la tradition des Humanités, ne signifie pas seulement commenter ses idées (la
théorie des Formes, l'allégorie de la caverne), mais comprendre comment ces idées émergent de la
langue même. C'est saisir la nuance entre eidos et idea, comprendre la valeur d'un aspect verbal, la
portée d'une négation.

Or, les conditions actuelles d'enseignement poussent vers une analyse purement conceptuelle ou
littéraire, détachée du texte source. On étudie le Phédon et l'immortalité de l'âme 12 à travers des
traductions, tout en prétendant faire du grec. L'élève apprend à disserter sur le « corps-tombeau »
(sôma-sèma) sans nécessairement être capable de repérer ce jeu de mots dans une phrase complexe ou
de comprendre la syntaxe qui l'environne. Cette déconnexion transforme l'enseignement du grec en un
enseignement de « philosophie sur documents traduits », tout en conservant l'étiquette prestigieuse
de « Langue et Culture de l'Antiquité ».



\subsection{La mécanique du simulacre : « sans savoir le lire »}

L'impasse n'est pas seulement un constat d'échec ; elle a généré une pédagogie de survie, un système
de compensation que l'on peut qualifier de « mécanique du simulacre ». Puisque l'on ne peut plus
lire, on fait « semblant », on utilise des béquilles, on remplace la langue par la culture.



\subsection{Du déchiffrage au « simulacre de lecture »}

La critique la plus acerbe porte sur la nature même de l'activité de lecture en classe. Faute d'une
exposition suffisante à la langue cible (input), les élèves ne développent pas d'automatismes de
lecture. Ils restent bloqués au stade du « déchiffrage ».\cite{I-3-3-ref13}

Le cours de grec devient alors un exercice de cryptographie : l'élève, armé d'un dictionnaire et
d'une grammaire, tente de résoudre une énigme mot par mot, cherchant le verbe, puis le sujet, sans
jamais construire de représentation mentale immédiate du sens. Cette activité, intellectuellement
exigeante mais linguistiquement stérile en termes de fluidité, est qualifiée de « simulacre de
lecture » par certains pédagogues.\cite{I-3-3-ref14} Elle mime l'acte de lire mais en ignore les
processus cognitifs réels (anticipation, inférence, compréhension globale).

Contrairement aux méthodes « nature » ou « directes » (comme la méthode Ørberg pour le latin ou
Athénaze pour le grec), qui prônent une immersion et une lecture cursive progressive 15,
l'enseignement français reste majoritairement attaché à la méthode « grammaire-traduction » 16,
inefficace avec des horaires réduits.



\subsection{L'institutionnalisation de la béquille : le texte bilingue}

L'aveu le plus flagrant de cette incapacité à lire est la généralisation de l'usage du texte
bilingue (texte grec et sa traduction en regard) dans les évaluations et les cours.

La réforme du concours d'entrée à l'École Normale Supérieure (ENS) en est l'illustration parfaite.
L'épreuve de langue et culture anciennes a été modifiée pour inclure un commentaire sur un texte
fourni « sous forme entièrement bilingue ».\cite{I-3-3-ref17} Si l'exercice de version subsiste pour
les spécialistes, cette nouvelle modalité pour le tronc commun des khâgneux reconnaît implicitement
que l'élite des étudiants littéraires n'est plus capable de commenter un texte grec de manière
autonome.

Le texte bilingue agit comme une « béquille ».\cite{I-3-3-ref19} L'élève ou l'étudiant lit la
traduction française, repère les mouvements du texte, puis retourne vers le grec pour piocher
quelques termes techniques ou connecteurs logiques afin de « faire grec » dans son commentaire. Il
ne lit pas le grec ; il vérifie la présence du grec dans la traduction. C'est un « semblant de
traduction » et d'analyse 20, une validation par procuration qui permet de sauver les apparences de
la compétence philologique.



\subsection{La « civilisation » comme refuge et cache-misère}

Face à l'aridité d'une langue devenue inaccessible, l'enseignement se replie massivement sur la «
civilisation » ou la « culture ». Les rapports de jury notent que chez les candidats fragiles, «
l'axe culturel prend souvent une place prépondérante au détriment de la langue ».\cite{I-3-3-ref10}

Cette dérive culturaliste transforme le cours de langue en un cours d'histoire ou d'anthropologie de
l'Antiquité dispensé en français. Si la culture est essentielle, elle sert ici d'alibi. On parle des
institutions athéniennes, de la mythologie, de la condition féminine chez Aristophane, en s'appuyant
sur des œuvres lues en traduction intégrale.

Florence Dupont et Barbara Cassin ont critiqué cette approche patrimoniale figée. Dans leur appel
pour une refondation des Humanités, elles dénoncent une Antiquité « conservatoire » et plaident pour
une Antiquité « laboratoire » qui accepterait le métissage et l'altérité.\cite{I-3-3-ref21}
Cependant, leur critique de l'enseignement traditionnel (jugé élitiste et obsédé par la grammaire) a
parfois été interprétée comme une validation de l'abandon de l'exigence linguistique stricte au
profit d'une approche plus « culturelle » et interdisciplinaire, ce qui, paradoxalement, peut
renforcer le recul de la maîtrise de la langue elle-même si les moyens horaires ne suivent
pas.\cite{I-3-3-ref21}



\subsection{L'exemple paradigmatique de platon : le sommet de l'illusion}

Pourquoi le point I.3.\cite{I-3-3-ref3} cite-t-il spécifiquement Platon? Parce que l'auteur de la
\textit{République} représente le point de tension maximal entre la perte des moyens et le maintien
des prétentions.



\subsection{L'inaccessibilité linguistique de platon}

Platon n'est pas seulement un philosophe ; c'est un prosateur attique d'une sophistication extrême.
Sa langue est un instrument de précision où la nuance philosophique repose entièrement sur la
syntaxe.

\begin{itemize}\item \textbf{La particule et la nuance :} Un dialogue platonicien est structuré par des particules (\textit{ge}, \textit{dè}, \textit{mèn oun}) qui indiquent l'ironie, l'assentiment forcé, la conclusion logique ou la concession. Sans une maîtrise avancée de ces micro-outils (qui demande des années de pratique), le lecteur passe à côté de la dynamique dialectique.\item Le jeu sur le lexique : Platon forge des concepts à partir de la langue courante. Comprendre l'Idée du Bien (hè tou agathou idea) nécessite de sentir la résonance visuelle du mot idea (ce qui est vu, l'aspect) pour saisir le glissement métaphysique.\end{itemize}Dans un enseignement qui a « perdu ses moyens », ces nuances sont invisibles. L'élève ne voit que
des mots grecs opaques qu'il traduit laborieusement par des concepts français figés (« Idée », «
Forme »), perdant tout le travail de la pensée en acte dans la langue.



\subsection{L'analyse du \textit{phédon} : une étude de cas}

Prenons l'exemple de l'analyse du Phédon sur l'immortalité de l'âme, souvent citée dans les
programmes universitaires ou de lycée.\cite{I-3-3-ref12}

L'objectif affiché est d'analyser l'argumentation de Socrate. L'élève va étudier la structure
logique des arguments (l'argument des contraires, l'argument de la réminiscence). Mais sans savoir
lire le grec, il ne peut pas analyser comment Platon construit cette immortalité grammaticalement.
Il ne verra pas l'usage des modes verbaux qui marquent le degré de certitude ou d'hypothèse. Il sera
condamné à répéter le commentaire du manuel ou de l'enseignant sur le texte français.

Ainsi, prétendre « analyser Platon » sans compétence de lecture autonome revient à commenter une
partition de musique sans savoir la déchiffrer, en se basant uniquement sur l'écoute d'un
enregistrement (la traduction). C'est une activité culturelle respectable, mais ce n'est pas de la
philologie, et encore moins de l'analyse de texte originale.



\subsection{Conséquences et perspectives : sortir du fossile?}

L'enseignement fossile décrit ici est à la croisée des chemins. Le maintien du statu quo — des
prétentions élevées avec des moyens dérisoires — condamne la discipline à une mort lente par
asphyxie et hypocrisie sociale.



\subsection{La sélection par l'échec et l'élitisme social}

En maintenant l'exigence de la version et de l'analyse de haut vol (Platon) sans donner les moyens
horaires de l'acquérir à l'école, le système renvoie la réussite à l'héritage culturel familial ou
aux établissements d'élite (souvent privés ou grands lycées parisiens) qui parviennent à maintenir
des conditions d'enseignement privilégiées. Le grec devient un marqueur de distinction sociale
d'autant plus efficace qu'il est officiellement « offert à tous » mais réellement accessible à très
peu. L'accusation d'élitisme, que réfutent les défenseurs des langues anciennes, est ainsi validée
par la structure même de l'offre éducative.\cite{I-3-3-ref3}



\subsection{Vers une refondation didactique ou une mort définitive?}

Pour sortir de l'impasse, des voix s'élèvent pour proposer des alternatives, inspirées notamment des
travaux de Pierre Judet de La Combe et Heinz Wismann sur « l'avenir des langues ».\cite{I-3-3-ref23}
Ils plaident pour ne pas considérer les langues anciennes comme un patrimoine sacré à conserver (le
fossile), mais comme un outil de décentrement et de formation de l'esprit critique.

Cependant, cela implique une refondation didactique :

\begin{itemize}\item \textbf{Adopter des méthodes inductives :} Abandonner le fétichisme de la grammaire-traduction au profit de méthodes favorisant la lecture cursive et l'imprégnation (type Ørberg ou \textit{Reading Greek}) 15, pour redonner une compétence de lecture réelle, même sur un périmètre lexical plus modeste.\item \textbf{Redéfinir les objectifs :} Accepter de réduire les prétentions (moins de version pure, moins d'auteurs canoniques inaccessibles) pour garantir que ce qui est appris est réellement maîtrisé. Mieux vaut savoir lire couramment un texte simple d'Ésope ou de Xénophon que de déchiffrer péniblement trois lignes de Platon avec une traduction.\item \textbf{L'approche « Laboratoire » :} Suivre la piste de Florence Dupont d'une Antiquité-laboratoire 22, où l'on travaille sur les écarts anthropologiques, mais en ancrant ce travail dans une pratique de la langue qui ne soit pas un simulacre.\end{itemize}

\subsection{Conclusion}

L'analyse du point I.3.\cite{I-3-3-ref3} révèle une vérité inconfortable : l'enseignement du grec
ancien en France fonctionne aujourd'hui largement sur un mensonge institutionnel. C'est un
enseignement fossile parce qu'il a gardé la coquille vide d'une ambition humaniste (lire les Grands
Textes) alors que le corps vivant de la discipline (le temps, la pratique, l'immersion) s'est
décomposé sous l'effet des réformes comptables et pédagogiques.

« Analyser Platon sans savoir le lire » n'est pas seulement un paradoxe ; c'est le symptôme d'un
système qui a renoncé à transmettre les moyens de sa propre reproduction, préférant la mise en scène
du savoir à la construction patiente de la compétence. Si l'on ne résout pas cette contradiction —
soit en rétablissant les moyens (improbable), soit en refondant les méthodes et les prétentions
(nécessaire) — le grec ancien cessera définitivement d'être une discipline formatrice pour devenir
une simple curiosité muséale, réservée à une poignée d'initiés et simulée pour le reste des élèves.

